% Part: first-order-logic
% Chapter: syntax-and-semantics
% Section: terms-formulas

\documentclass[../../../include/open-logic-section]{subfiles}

\begin{document}

\olfileid[fr]{fol}{syn}{frm}

\olsection{Termes et \printtoken{p}{formula}}

Une fois donné un langage du premier ordre~$\Lang L$, nous pouvons
définir des expressions construites à partir du vocabulaire de~$\Lang L$.
Parmi elles figurent notamment les \emph{termes} et les
\emph{!!{formula}s}.

\begin{defn}[Termes]
\ollabel{defn:terms}
L'ensemble des \emph{termes}~$\Trm[L]$ de~$\Lang L$ est défini
inductivement par les clauses suivantes :
\begin{enumerate}
\item Toute !!{variable} est un terme.
\item Toute !!{constant} de~$\Lang L$ est un terme.
\item Si $f$ est un !!{function} à $n$ places et si $t_1$, \dots,
  $t_n$ sont des termes, alors $\Atom{f}{t_1, \ldots, t_n}$ est un
  terme.
\tagitem{limitClause}{Rien d'autre n'est un terme.}{}
\end{enumerate}
Un terme ne contenant aucune !!{variable} est un \emph{terme clos}.
\end{defn}

\begin{explain}
Nous avons présenté les !!{constant}s comme une catégorie de symboles
distincte, tant dans la spécification du langage que dans la
définition des termes. Nous aurions aussi pu les traiter comme des
!!{function}s à zéro place ; la deuxième clause de la définition
des termes deviendrait alors superflue. Il suffirait de comprendre
$\Atom{f}{t_1, \ldots, t_n}$ comme le seul symbole~$f$ lorsque
$n = 0$.
\end{explain}

\begin{defn}[Formules]
\ollabel{defn:formulas}
L'ensemble des \emph{!!{formula}s}~$\Frm[L]$ du langage~$\Lang L$
est défini inductivement comme suit :
\begin{enumerate}
\tagitem{prvFalse}{$\lfalse$ est une !!{formula} atomique.}{}

\tagitem{prvTrue}{$\ltrue$ est une !!{formula} atomique.}{}

\item Si $R$ est un !!{predicate} à $n$ places de~$\Lang L$ et si
  $t_1$, \dots, $t_n$ sont des termes de~$\Lang L$, alors
  $\Atom{R}{t_1,\ldots, t_n}$ est une !!{formula} atomique.

\item Si $t_1$ et $t_2$ sont des termes de~$\Lang L$, alors
  $\Atom{\eq}{t_1, t_2}$ est une !!{formula} atomique.

\tagitem{prvNot}{Si $!A$ est !!a{formula}, alors $\lnot !A$ est
  !!a{formula}.}{}

\tagitem{prvAnd}{Si $!A$ et $!B$ sont des !!{formula}s, alors
  $(!A \land !B)$ est !!a{formula}.}{}

\tagitem{prvOr}{Si $!A$ et $!B$ sont des !!{formula}s, alors
  $(!A \lor !B)$ est !!a{formula}.}{}

\tagitem{prvIf}{Si $!A$ et $!B$ sont des !!{formula}s, alors
  $(!A \lif !B)$ est !!a{formula}.}{}

\tagitem{prvIff}{Si $!A$ et $!B$ sont des !!{formula}s, alors
  $(!A \liff !B)$ est !!a{formula}.}{}

\tagitem{prvAll}{Si $!A$ est !!a{formula} et $x$ !!a{variable},
  alors $\lforall[x][!A]$ est !!a{formula}.}{}

\tagitem{prvEx}{Si $!A$ est !!a{formula} et $x$ !!a{variable},
  alors $\lexists[x][!A]$ est !!a{formula}.}{}

\tagitem{limitClause}{Rien d'autre n'est !!a{formula}.}{}
\end{enumerate}
\end{defn}

\begin{explain}
Les définitions de l'ensemble des termes et de celui des !!{formula}s
sont des \emph{définitions inductives}. Pour construire l'ensemble
des !!{formula}s, nous procédons, en substance, par une infinité
d'étapes. À la première, nous posons que toutes les !!{formula}s
atomiques sont des !!{formula}s. Cela correspond aux premiers cas de
la définition, à savoir ceux de
\iftag{prvTrue}{$\ltrue$, }{}%
\iftag{prvFalse}{$\lfalse$, }{}%
$\Atom{R}{t_1,\dots,t_n}$ et $\Atom{\eq}{t_1,t_2}$. Par
« !!{formula} atomique », nous entendons donc toute !!{formula}
de l'une de ces formes.

Les autres clauses de la définition donnent des règles pour construire
de nouvelles !!{formula}s à partir de !!{formula}s déjà obtenues. À
la deuxième étape, elles permettent d'en former à partir des
!!{formula}s atomiques. À la troisième, on en forme à partir des
!!{formula}s atomiques et de celles obtenues à la deuxième étape, et
ainsi de suite. Une !!{formula} est une expression construite à l'une
de ces étapes, et rien d'autre.
\end{explain}

Par convention, nous écrivons $\eq$ entre ses arguments et supprimons
les parenthèses : $\eq[t_1][t_2]$ abrège
$\Atom{\eq}{t_1,t_2}$. De même, $\lnot \Atom{\eq}{t_1,t_2}$
est abrégé par $\eq/[t_1][t_2]$. Lorsque la !!{formula}
$(!B \ast !C)$ est construite à partir de $!B$ et de $!C$ au moyen
d'un connecteur binaire~$\ast$, nous omettrons souvent la paire de
parenthèses extérieure et écrirons simplement~$!B \ast !C$.

\begin{intro}
Certains manuels de logique exigent que la !!{variable}~$x$ apparaisse
dans~$!A$ pour que
\iftag{prvEx}{$\lexists[x][!A]$ }{}%
\iftag{notprvEx,notprvAll}{}{et }%
\iftag{prvAll}{$\lforall[x][!A]$ }{}%
\iftag{notprvEx,notprvAll}{soit considérée comme !!a{formula}}{soient
  considérées comme des !!{formula}s}. Cette exigence n'est pas
nécessaire et son absence simplifie les choses.
\end{intro}

\begin{tagblock}{defNot,defOr,defAnd,defIf,defIff,defTrue,defFalse,defEx,defAll}
\begin{defn}
Les !!{formula}s construites à l'aide de symboles définis se
comprennent de la manière suivante :

\begin{tagenumerate}{defTrue,defFalse,defNot,defOr,defAnd,defIf,defIff,defEx,defAll}

\tagitem{defTrue}{$\ltrue$ abrège
  \iftag{prvFalse}{$\lnot\lfalse$}{$(!A \lor \lnot !A)$ pour une
    !!{formula} atomique fixée~$!A$}.}{}

\tagitem{defFalse}{$\lfalse$ abrège
  \iftag{prvTrue}{$\lnot\ltrue$}{$(!A \land \lnot !A)$ pour une
    !!{formula} atomique fixée~$!A$}.}{}

\tagitem{defNot}{$\lnot !A$ abrège $!A \lif \lfalse$.}{}

\tagitem{defOr}{$!A \lor !B$ abrège
  \iftag{prvAnd}{$\lnot(\lnot !A \land \lnot !B)$}{$\lnot !A \lif
    !B$}.}{}

\tagitem{defAnd}{$!A \land !B$ abrège
  \iftag{prvOr}{$\lnot(\lnot !A \lor \lnot !B)$}{$\lnot (!A \lif
    \lnot !B)$}.}{}

\tagitem{defIf}{$!A \lif !B$ abrège
  \iftag{prvOr}{$\lnot !A \lor !B$}{$\lnot (!A \land \lnot !B)$}.}{}

\tagitem{defIff}{$!A \liff !B$ abrège $(!A \lif !B) \land (!B
  \lif !A)$.}{}

\tagitem{defAll}{$\lforall[x][!A]$ abrège $\lnot\lexists[x][\lnot !A]$.}{}

\tagitem{defEx}{$\lexists[x][!A]$ abrège $\lnot\lforall[x][\lnot !A]$.}{}
\end{tagenumerate}
\end{defn}
\end{tagblock}

Lorsque nous travaillons dans un langage destiné à une application
particulière, nous écrivons souvent les !!{predicate}s et les
!!{function}s à deux places entre leurs termes respectifs. C'est le
cas de $t_1 < t_2$ et de $(t_1 + t_2)$ dans le langage de
l'arithmétique, et de $t_1 \in t_2$ dans celui de la théorie des
ensembles. Par convention, le symbole de la fonction successeur en
arithmétique s'écrit même \emph{après} son argument :~$t'$. D'un
point de vue formel, il ne s'agit pourtant que d'abréviations
conventionnelles pour $\Atom{\Obj A^2_0}{t_1, t_2}$,
$\Obj f^2_0(t_1, t_2)$, $\Atom{\Obj A^2_0}{t_1, t_2}$ et
$f^1_0(t)$, respectivement.

\begin{defn}[Identité syntaxique]
Le symbole~$\ident$ exprime l'identité syntaxique de deux chaînes de
symboles : $!A \ident !B$ si et seulement si $!A$ et $!B$ ont la même
longueur et contiennent le même symbole à chaque position.
\end{defn}

De part et d'autre de~$\ident$, on peut placer des chaînes obtenues
par concaténation. Par exemple, $!A \ident (!B \lor !C)$ signifie
que la chaîne de symboles~$!A$ est identique à celle obtenue en
concaténant, dans cet ordre, une parenthèse ouvrante, la chaîne~$!B$,
le symbole~$\lor$, la chaîne~$!C$ et une parenthèse fermante. Si
c'est le cas, nous savons que le premier symbole de~$!A$ est une
parenthèse ouvrante, que $!A$ contient $!B$ comme sous-chaîne à
partir de son deuxième symbole, que cette sous-chaîne est suivie
de~$\lor$, et ainsi de suite.

Les termes et les !!{formula}s étant construits à partir d'éléments
de base par des définitions inductives, nous pouvons appliquer les
principes d'induction suivants pour démontrer leurs propriétés.

\begin{lem}[\emph{Principe d'induction sur les termes}]
    \ollabel{lem:trmind}
    Soit $\Lang L$ un langage du premier ordre.
    Si une propriété~$P$ vérifie les conditions suivantes :
    %
    \begin{enumerate}
        \item elle vaut pour toute !!{variable}~$v$ ;
        %
        \item elle vaut pour toute !!{constant}~$a$ de~$\Lang L$ ;
        %
        \item elle vaut pour $f(t_1,\dotsc,t_n)$ dès qu'elle vaut
        pour $t_1$,~\dots, $t_n$ et que $f$ est un !!{function}
        à $n$ places de~$\Lang L$,
    \end{enumerate}
    en supposant que $t_1$,~\dots, $t_n$ soient des termes de~$\Lang{L}$,
    alors $P$ vaut pour tout terme de~$\Trm[L]$.
\end{lem}

\begin{prob}
    Démontrer le \olref[fol][syn][frm]{lem:trmind}.
\end{prob}

\begin{lem}[\emph{Principe d'induction sur les !!{formula}s}]
    \ollabel{thm:frmind}
    Soit $\Lang L$ un langage du premier ordre.
    Si une propriété~$P$ vaut pour toutes les !!{formula}s atomiques
    et vérifie les conditions suivantes :
    %
    \begin{enumerate}
        \tagitem{prvNot}{elle vaut pour $\lnot !A$ dès qu'elle
            vaut pour~$!A$ ;}{}
        \tagitem{prvAnd}{elle vaut pour $(!A \land !B)$ dès qu'elle
            vaut pour $!A$ et~$!B$ ;}{}
        \tagitem{prvOr}{elle vaut pour $(!A \lor !B)$ dès qu'elle
            vaut pour $!A$ et~$!B$ ;}{}
        \tagitem{prvIf}{elle vaut pour $(!A \lif !B)$ dès qu'elle
            vaut pour $!A$ et~$!B$ ;}{}
        \tagitem{prvIff}{elle vaut pour $(!A \liff !B)$ dès qu'elle
            vaut pour $!A$ et~$!B$ ;}{}
        \tagitem{prvEx}{elle vaut pour $\lexists[x][!A]$ dès qu'elle
            vaut pour~$!A$ ;}{}
        \tagitem{prvAll}{elle vaut pour $\lforall[x][!A]$ dès qu'elle
            vaut pour~$!A$ ;}{}
  \end{enumerate}
  en supposant que $!A$ et $!B$ soient des !!{formula}s de~$\Lang{L}$,
  alors $P$ vaut pour toutes les !!{formula}s de~$\Frm[L]$.
\end{lem}

\end{document}
